\begin{trapbox}

	\begin{description}[style=nextline, leftmargin=2.8em, labelsep=0.5em, itemsep=0.35em]
		\item[\textbf{\textcolor{CTrap}{易错点一（忽略定义域）}}]
			在 $x=\pi+2k\pi$ 时使用万能公式，此时 $\tan\frac{x}{2}$ 无意义，却仍代入 $t$ 计算。
		\item[\textbf{\textcolor{CTrap}{正确理解（定义域先行）：}}]
			使用前必须检查 $x\neq \pi+2k\pi$，确保 $t$ 有定义。
		\item[\textbf{\textcolor{CTrap}{错因分析（机械套用）：}}]
			误认为万能公式是恒等式，忽视了半角正切的存在条件。
	\end{description}

	\medskip
	\begin{description}[style=nextline, leftmargin=2.8em, labelsep=0.5em, itemsep=0.35em]
		\item[\textbf{\textcolor{CTrap}{易错点二（遗漏正切条件）}}]
			使用 $\tan x = \frac{2t}{1-t^2}$ 时没有注意 $x\neq \frac{\pi}{2}+k\pi$，当 $t^2=1$ 时分母为零，公式失效。
		\item[\textbf{\textcolor{CTrap}{正确理解（分母检查）：}}]
			使用正切公式前，应验证 $x\neq \frac{\pi}{2}+k\pi$ 或 $t^2\neq 1$。
		\item[\textbf{\textcolor{CTrap}{错因分析（条件遗忘）：}}]
			只记住公式形式，忘记了推导中 $\cos x$ 不能为零的条件。
	\end{description}

	\medskip
	\begin{description}[style=nextline, leftmargin=2.8em, labelsep=0.5em, itemsep=0.35em]
		\item[\textbf{\textcolor{CTrap}{易错点三（回代误解周期）}}]
			从 $t$ 值反求 $x$ 时，直接写 $x=2\arctan t$ 而忽略周期性，遗漏 $+2k\pi$。
		\item[\textbf{\textcolor{CTrap}{正确理解（通解形式）：}}]
			应写 $x=2\arctan t + 2k\pi$ ($k\in \mathbb{Z}$)，若题目指定范围再筛选。
		\item[\textbf{\textcolor{CTrap}{错因分析（周期忽视）：}}]
			默认反正切的主值范围，忘记正切函数的周期为 $\pi$，从而 $x$ 周期为 $2\pi$。
	\end{description}

\end{trapbox}
